\section{Wilson 线、穿刺点与 Borel--Weil--Bott}
\label{sec:wilson}

前两节只处理了“没有纽结”的理论。要让纽结进来，就得在三维流形里放 Wilson 线。这一节回答三个问题：Wilson 线怎么变成曲面上一个穿刺点？穿刺点上的约束为什么不能经典求解？Witten 用什么办法把它“反量子化”回去？

\subsection{纽结变成一个穿刺点}

把 $M$ 沿曲面 $\Sigma$ 切开时，一条穿过 $\Sigma$ 的 Wilson 线会留下一个交点 $P$。在交点附近，规范场的行为由所在表示 $R$ 决定：横向分量与 $R$ 的李代数生成元耦合。所以“带 Wilson 线的三维理论”等价于“带标记点的曲面理论”，每个标记点挂一个表示 \src{371}。

标记点上的高斯定律不再是 $F=0$，而是带源的形式 \src{371}（(3.4)）：
\begin{equation}
\frac{k}{8\pi}\,\epsilon^{ij}F^a_{ij}(x)=\sum_s \delta^2(x-P_s)\,T^a_{(s)} ,
\label{eq:markedgauss}
\end{equation}
其中 $T^a_{(s)}$ 是第 $s$ 个标记点上表示 $R_s$ 的生成元。

\subsection{为什么不能经典求解}

约束右端是\textbf{算符值}：$T^a_{(s)}$ 之间不对易，它作用在表示空间上。所以 (3.4) 不是一个可以“解出某个经典联络”的方程，而是一个量子约束 \src{372}。原文的说法是：若硬把 $A$ 当作经典对象去解，得到的不是普通联络，而是“量子群”的某种东西；但 Witten 明说这条路只是提出来，没有构造 \src{372}。

第二个困难是奇异性：点源使 $A$ 在 $P_s$ 处奇异，平坦性只在 $\Sigma\setminus\{P_s\}$ 上成立。所以相空间不是“$\Sigma$ 上平坦联络模空间”，而是“带奇点的平坦联络模空间”。

\subsection{Witten 的办法：把表示也量子化}

Witten 的解法不是解约束，而是换一个角度看表示：\textbf{把 $R$ 本身看成某个辛流形的几何量子化} \src{371--372}。

具体地，取旗流形 $G/T$（$T$ 为极大环面）。它有自然的 K\"ahler 结构；对支配整权 $\lambda$，取辛形式 $\omega_\lambda$，其 de Rham 类就是 $\lambda$。Borel--Weil--Bott 定理说
\begin{equation}
H^0(G/T,\cL_\lambda)=R_\lambda ,
\end{equation}
即 $\omega_\lambda$ 的几何量子化恰好给出最高权为 $\lambda$ 的不可约表示。注意这里 $\lambda$ \textbf{不加 $\rho$}：$\rho$ 的位移属于半形式修正那一层（$K_{G/T}^{1/2}=\cL_{-\rho}$），不在指标里。

于是可以构造一个更大的相空间 $M_{P_i,R_i}$：在 $\Sigma\setminus\{P_i\}$ 上取平坦 $G$ 联络，在每个 $P_i$ 处把结构群约化到 $T$，参数由 $\lambda_i$ 给出。这个空间上的经典约束 (3.4) 有意义；量子化之后，得到的恰好是 Segal 的流代数共形块空间 \src{372--373}。

\begin{keypoint}
\textbf{这一步的逻辑}：不是在三维理论里“解”穿刺约束，而是把“表示”从外部数据变成相空间的一部分。原来 Wilson 线 $W_R(C)=\Tr_R P\exp\int_C A$ 里的 $\Tr_R$ 与 $R$，现在由 $G/T$ 这一层的量子化提供。Witten 在 p.371 说这是 “should” 的做法，并承认 “deferring a fuller treatment”——所以这是\textbf{纲领 + 部分实现}，不是完整构造。
\end{keypoint}

\subsection{可积表示条件}

只有“可积”的表示才有非零截面：$SU(2)$ 时 $2j\le k$，一般地 $\lambda$ 属于 level $k$ 的仿射 Weyl 群基本域。超出这个范围的表示在共形块里“退耦”（对应维数为零）\src{372}。

\textbf{这个条件不是 Borel--Weil--Bott 的推论}——BWB 对任意支配整权都成立。它来自仿射李代数：level $k$ 的可积表示要求 $(\lambda,\theta)\le k$（河野命题 1.5，书 p.11）。原文把它写在 p.372 的括号里，一笔带过；Moore 中译 p.120 与上野 p.39 才明写。

\subsection{记号对照}

读原文时容易把几个“$\lambda$”混起来：
\begin{center}
\small
\begin{tabular}{@{}lll@{}}
\toprule
对象 & 记号 & 说明 \\
\midrule
最高权 & $\lambda$ & 表示 $R_\lambda$ 的最高权；BWB 里不加 $\rho$ \\
辛形式 & $\omega_\lambda$ & $G/T$ 上，$[\omega_\lambda]=\lambda$ \\
共形块 & $H^0(\cM_J,\cL^{\otimes k})$ & 带标记点时是子空间，由 $\lambda_i$ 决定 \\
仿射 level & $k+h^\vee$ & 半形式修正后；本文多处写 “level $k$” 是简写 \\
\bottomrule
\end{tabular}
\end{center}

\begin{boundary}
\begin{itemize}[leftmargin=1.4em]
\item 原文给出：(3.4)、BWB 实现、扩大相空间、引 Segal、可积性括号。
\item 我们补的：把“$\lambda$ 不加 $\rho$”与半形式修正分开（$\rho$ 属于 $K^{1/2}$）；把 $2j\le k$ 归因于仿射代数而非 BWB。
\item 未证：把 Wilson 线换成“$S\to G/T$ 的映射的泛函积分”这一纲领，原文只有 “should” 与 “deferring a fuller treatment”，没有构造；“量子群”评论同样没有构造。
\end{itemize}
\end{boundary}
